\documentclass[]{acoust}
\usepackage{url}
\usepackage{booktabs}
%%% [letter] for Acoustical Letter
%%% [shortnote] for Short Note
%%% [techrep] for Technical Report
%%% Otherwise, Paper
%%% [prepr] for onecolumn layout = Old format for submission to AST
\setlength\tabcolsep{4pt} % add more space between columns of table

\def\header{Minimizing data-acquisition for USCT imaging} % A short running title of not more than 50 keystrokes

% The first letter of the title and all proper nouns should be capitalized with the rest in lower case.
% As a general rule, abbreviations should not be used in the Title.
\title{Towards minimizing data-acquisition for a pixel-based ultrasound computed tomography echo imaging method}
% \subtitle{Sub-title goes here}

\author{Tianhan Tang$^{1,2}$\thanks{tianhantang.pd@gmail.com}, Ichiro Sakuma$^{1}$, Takashi Azuma$^{2}$, Shin-Ichiro Umemura$^{2}$,\\
Atsushi Otsubo$^{2}$, Toshihide Iwahashi$^{2}$, Hirofumi Nakamura$^{2}$}

\inst{$^{1}$Department of Bioengineering, Graduate School of Engineering, The University of Tokyo\\
Room 722, Engineering 14th Bldg., 7-3-1 Hongo, Bunkyo-ku, Tokyo, 113-8656 Japan\\
$^{2}$Lily MedTech Inc.\\
Entrepreneur Lab, South Research Bldg., 7-3-1 Hongo, Bunkyo-ku, Tokyo, 113-8485 Japan
}

% \author{% List of author goes here.
% Author1$^1$, Author2$^2$, and Author3$^1$\thanks{E-mail address(es) must be a footnote.}}
% 
% \inst{% List of affiliation(s) goes here. \\
% $^{1}$ Multiple affiliation(s) must be numbered with superscript \\
% $^{2}$ Affiliation(2)}

%\recdate{}{}{} %% received date with yyyymmdd or yyyymd
%\accdate{}{}{} %% received date with yyyymmdd or yyyymd
%\jstdate{}{}{} %% J-STAGE Advance published date with yyyymmdd or yyyymd

\abst{
% >>> abstract
@anchor:t01
% <<<
}

\kword{Ultrasound computed tomography, Data reduction, System optimization, Inverse problem, Synthetic aperture}

%%%% The following three commands must be input only when a covering letter 
%%%% is required.

\mladdress{Room 722, Engineering 14th Bldg., 7-3-1 Hongo, Bunkyo-ku, Tokyo, 113-8656 Japan} %% Japanese is available.

\class{10} %%% A number from 1 to 7 should be typed for the appropriate classification.
         %%% 1 for Ultrasonic
         %%% 2 for Underwater Acoustics
         %%% 3 for Speech
         %%% 4 for Hearing
         %%% 5 for Speech and Hearing
         %%% 6 for Musical Acoustics
         %%% 7 for Electro-acoustics
         %%% 8 for Architectural Acoustics
         %%% 9 for Noise and Vibration
         %%% 10 for Acoustic Imaging
         %%% 11 for Education in Acoustics
         %%% 12 for Acoustical Assists(Acoustic Barrier-Free)
         %%% 13 for Thermoacoustic Technology
         %%% 14 for Sound Design

\recommend{2}  %%% If this paper is recommended as an Acoustical Letter by the Editorial Board, 
               %%% input "1"; otherwise, input "2".  This command is valid for Acoustical Letter
			   %%% only.
%%%%%%%%%%%%%%%%%%%%%%%%%%%%%%%%%%%%%%%%%%%%%%%%%%%%%%%%%%%%%%%%%%% 

\begin{document}
\maketitle

\section{Introduction}
\label{sec:introduction}

% >>> introduction
@anchor:t02
% <<<

\section{Method}
\label{sec:method}

\subsection{The pixel-based echo imaging method}
% >>> method/p1
@anchor:t03
% <<<

\subsection{Minimizing data-acquisition}
% >>> method/p2
@anchor:t04
% <<<

\subsection{The simulation experiment}
% >>> method/p3
@anchor:t05
% <<<

\section{Result}
\label{sec:result}

\subsection{Condition numbers of systems with different configurations}
% >>> result/p1
@anchor:t06
% <<<

\subsection{Phantom data reconstruction under different SNR}
% >>> result/p2
@anchor:t07
% <<<

\subsection{Comparison between pixel-based imaging method and synthetic aperture method}
% >>> result/p3
@anchor:t08
% <<<

\section{Discussion}
\label{sec:discussion}

\subsection{Advantages and features of the pixel-based imaging method}
% >>> discussion/p1
@anchor:t09
% <<<

\subsection{Challenges and future directions for the pixel-based imaging method}
% >>> discussion/p2
@anchor:t10
% <<<

\section{Conclusion}
\label{sec:conclusion}

% >>> conclusion
@anchor:t12
% <<<

%% If the number of items less than 10, "9" must be input
%% Less than 100, input "99".
\begin{thebibliography}{99} 

% >>> reference
@anchor:t13
% <<<
\end{thebibliography}

% \covering %%% This command will produce a covering letter.

\end{document}
